% ============================================================================
\clearpage
\section*{v4.10.0-r10 Source-Only: The Bridge of the Counting Principle, and the Counting Rule as a Frame-Dependent Area Quantum}
\addcontentsline{toc}{section}{v4.10.0-r10 Source-Only: The Bridge of the Counting Principle, and the Counting Rule as a Frame-Dependent Area Quantum}
\label{sec:v4100-r10-area-quantum}
% ============================================================================

\paragraph{Purpose.}
Remark v4.10.0-r09.2 left one bridge open: the principle P lets the tachyon frame count only the smallest area excitation because it sits at its
floor (v4.9.0-r63.3), and links floor energy to horizon energy through the local first law.  This layer examines that bridge.  It does not hold in
that form; what remains is a sharper statement of the counting rule.  A parallel audit on GitHub reached the same verdict (r10.4).  Exploration, analytic, no criterion.  The notes are
\texttt{notes/\allowbreak research/\allowbreak v4\_10\_0\_r10/}.  Earlier layers are not rewritten.

\noindent\textbf{Source-only status.}
This v4.10.0-r10 layer is source-only.  It does not build a PDF, does not assign a new DOI,
and does not authorize a public release.  The v4.10.0 Version DOI
\href{https://doi.org/10.5281/zenodo.23991357}{10.5281/zenodo.23991357} and the Concept DOI
\href{https://doi.org/10.5281/zenodo.19210728}{10.5281/zenodo.19210728} are unchanged by it.

\subsection*{v4.10.0-r10 results}
\label{subsec:v4100-r10-results}

\paragraph{Theoremlet v4.10.0-r10.1 (the counting equation is an Unruh statistics, and the distance cancels).}
Analytic.  By the local first law (Frodden, Ghosh and Perez, 2013) an observer at proper distance \(\ell\) from the horizon assigns a puncture of spin
\(j\), area \(8\pi\gamma\,a(j)\) (\(\ell_P=1\)), the energy \(E_j=\gamma\,a(j)/\ell\), and has the Unruh temperature \(T_U=1/(2\pi\ell)\).  Hence
\[
E_j/T_U=2\pi\gamma\,a(j),
\]
independent of \(\ell\) (checked for \(\ell=0.5\) to \(1000\)).  The counting equation \(\sum_j g(j)\,e^{-2\pi\gamma a(j)}=1\) is this statistics (Ghosh and
Perez, 2011).  All static observers count alike: a higher or lower energy shifts \(E_j\) and \(T_U\) in the same ratio.

\paragraph{Theoremlet v4.10.0-r10.2 (the tachyon frame sees the same statistics).}
By v4.10.0-r07.2 the tachyon map commutes with the boost and fixes the vacuum; an observer in frame T sees the same vacuum with the same modular group,
hence the same KMS temperature.  The floor of r63.3 (particle energy \(E\to0\) at fixed invariant) is a property of a tachyon particle, not of the state an
observer sees at the horizon.  \textbf{Correction to Remark r09.2:} the bridge in P2 does not hold through the local first law and the Unruh temperature;
along that route M and T count alike.  ``T reaches only the smallest excitation'' would need a state of T that is not thermal with respect to the boost;
the vacuum is thermal.

\paragraph{Theoremlet v4.10.0-r10.3 (the counting rule as a frame-dependent area quantum).}
Analytic.  If both frames see the same statistics (r10.1, r10.2), the same geometry (r07.2) and both satisfy \(S=A/4\) (r06.2), the only remaining place for a
difference is the area quantum.  The assignment of counting rules to the frames (r08.3, r09) is then equivalent to a frame-dependent Barbero--Immirzi
parameter,
\[
\gamma_M=\gamma_{\rm BI},\qquad \gamma_T=\gamma_{\rm string},\qquad \gamma_M/\gamma_T=\gamma_L=1.864700:
\]
the smallest area quantum is smaller by the factor \(\gamma_L\) in frame T than in frame M, while the classical area and \(G\) are the same.  This is
v4.9.0-r64.1(b) (``the frame difference sits in \(\gamma\)'') in sharper form.  The duality rotation alone cannot change an area quantum, since it fixes
everything geometric; a reason would have to come from the boost part of the frame map (T as the floor of the mirror frame, r63.4), which only
reparametrizes the horizon.  No such reason is known.

\paragraph{Remark v4.10.0-r10.4 (relation to the parallel audit on GitHub).}
A source reconstruction of OP5 made with another language model (notes \texttt{op5\_\allowbreak spin\_\allowbreak selection\_\allowbreak audit\_\allowbreak 2026\_\allowbreak 11\_\allowbreak 06.md},
\texttt{op5\_\allowbreak common\_\allowbreak state\_\allowbreak access\_\allowbreak 2026\_\allowbreak 11\_\allowbreak 06.md}, \texttt{op5\_\allowbreak projector\_\allowbreak spin\_\allowbreak bridge\_\allowbreak 2026\_\allowbreak 11\_\allowbreak 06.md} in
\texttt{notes/\allowbreak research/}) reached \texttt{main} after layer r09.  It agrees with r10.2 and adds two points that bear on r09.2.
\begin{enumerate}[leftmargin=2em]
\item[(a)] Uncertainty saturation alone does not select \(j=\tfrac12\): \(|1,1\rangle\) and every \(|j,\pm j\rangle\) saturate as well.  ``Absolute minimum over \(j\)'' for T and
``extremal states at every \(j\)'' for M are two different rules; saturation is compatible with both and decides between neither.  P1 in r09.2 was used
only to exclude the weight \(2j+1\), which stands; the selection of \(j=\tfrac12\) for T rested on P2.
\item[(b)] The tachyon floor \(E=\mu/\sqrt{v^2-1}\to0\) does not contain \(j\); without a law linking \(\mu\) to the spin and a transition that changes \(j\), it gives no
relaxation to \(j=\tfrac12\).  This is the particle side of r10.2.
\item[(c)] A passive, invertible change of description cannot remove spin sectors; a T reading restricted to \(j=\tfrac12\) is either a restricted readout of the
same state (the rest stays in the balance) or a selection with its own weight.
\end{enumerate}
Together: P2 holds neither as an energy condition (r10.2) nor as a relaxation (b); P1 does not select the spin (a).  The singular position of r09.1 (H98)
is not affected; the assignment remains an assumption, now stated as a frame-dependent area quantum (r10.3).

\paragraph{Proof and measurement sources.}
\texttt{Erkundung\_\allowbreak Bruecke\_\allowbreak P2\_\allowbreak Unruh.md}; the cancellation of \(\ell\) in r10.1 is one line and was checked numerically.  \(\square\)

\paragraph{Consequence.}
Principle P keeps P1 (uncertainty saturation) as a reading; P2 as an energy condition is not supported along the route of the first law and should be
replaced by the plainer statement that each frame has its own area quantum.  That is an assumption, not a derivation.  The singular position found in
H98 is unaffected.  The three predictions on the \(\gamma_L\) core (H15, H46, H96) now visibly rest on one assumption: a frame-dependent area quantum with
ratio \(\gamma_L\).

\subsection*{v4.10.0-r10 status in the three frames}
\label{subsec:v4100-r10-frames}

\begin{itemize}[leftmargin=2em]
\item M and T: the same Unruh statistics and the same modular group (r10.1, r10.2); this part is frame free.
\item The area quantum differs: \(\gamma_M=\gamma_{\rm BI}\), \(\gamma_T=\gamma_{\rm string}\) (r10.3, assumption).
\item P: carries no counting.
\end{itemize}

\subsection*{v4.10.0-r10 register entries}
\label{subsec:v4100-r10-register}

\begin{center}
\small
\begin{tabular}{@{}>{\raggedright\arraybackslash}p{0.34\textwidth}>{\raggedright\arraybackslash}p{0.58\textwidth}@{}}
\toprule
statement & status after v4.10.0-r10\\
\midrule
bridge in P2 via first law and Unruh temperature & does not hold (correction to r09.2)\\
counting rule per frame & equivalent to a frame-dependent area quantum, \(\gamma_M/\gamma_T=\gamma_L\) (assumption)\\
reason for a frame-dependent area quantum & open\\
\bottomrule
\end{tabular}
\end{center}

\subsection*{v4.10.0-r10 ledger and guard update}
\label{subsec:v4100-r10-ledger-guard}

\[
R_{\rm chk}^{v4100\text{-}r10}=(1,1,0)
\]
for (bridge examined; counting rule restated; its reason open).

The layer spends no guard:
\[
R_{\rm glob}^{(v4100\text{-}r10,+)}=(0,1,1,1,1),
\qquad
\Gamma_{\rm guard}^{v4100\text{-}r10}=(1,1,1,1,1),
\qquad
\beta_4=1/256\ \text{unspent}.
\]
It does not imply exact finite \(C5\) zero closure, exact all-readout terminality,
\(C6/T2\), a global photon-clock/RH theorem, or global all-readout closure.
